\chapter{Introduction}
\section{Motivation}
Survival analysis studies the time until an event of interest, such as death,
disease recurrence, or equipment failure. In many studies, observation does not
begin at the origin of the relevant time scale. An individual may enter a study
several years after diagnosis while follow-up is measured from diagnosis.
Such an individual can appear in the study only if the event has not occurred
before entry. This observation mechanism is called left truncation, or delayed
entry.

Left truncation differs from right censoring. With right censoring, an
individual is observed but the event time is known only to exceed the end of
follow-up. With left truncation, individuals whose event precedes entry are
absent from the observed sample. Accounting for delayed entry therefore
requires attention both to sample selection and to who is under observation
at each event time.

For example, consider two individuals with observed follow-up ending at years
three and six. The second individual enters at year four. At an event at year
three, the second individual should not contribute to the risk set: they were
not yet under observation. An analysis using only the end of follow-up would
incorrectly count both individuals as at risk. This changes the denominators
used to estimate hazard increments and can change the splits chosen by a
survival tree.

Random survival forests combine survival trees to estimate outcomes from
covariates. The original formulation addresses right-censored outcomes
\citep{ishwaran2008}. Tree methods have also been developed specifically for
left-truncated, right-censored data \citep{fu2017}. This thesis focuses on
bringing delayed-entry support into the main fitting interface and
computational core of \pkg{}. The contribution is an applied software
extension and its evaluation, rather than a claim to introduce the first
tree-based method for left truncation.

\section{Objectives and research questions}
The primary objective is to extend the package so that the default survival
forest uses subject-specific entry times in tree construction and terminal-node
estimation. A second objective is to establish how this correction affects
prediction under controlled simulation conditions. A third is to document a
reproducible implementation suitable for a future upstream contribution.

The study addresses three questions:
\begin{enumerate}
\item Can delayed entry be incorporated into the existing R and C architecture
while preserving the ordinary fitting path when all entry times are zero?
\item How does supplying the true entry times affect estimated survival curves
and risk discrimination under different data-generating mechanisms?
\item Which package features and evaluation procedures require additional work
before delayed-entry support can be considered complete?
\end{enumerate}

\section{Scope and organization}
The current implementation concerns single-event survival outcomes and the
default log-rank splitting rule. Competing risks, time-varying covariates, and
alternative splitting rules are outside the validated scope of this draft.
Upstream submission is a future objective and is not reported as completed.

Chapter 2 introduces the statistical setting and related work. Chapter 3
describes the method and its implementation. Chapter 4 specifies the simulation
study, Chapter 5 presents the completed 100-replication results, and Chapter 6 discusses
limitations and remaining work.
